% Pursuing Stacks (cotes 134-2 to 134-8): archive cote 134-5, site page 53 (PDF page 54), typescript page 240.
% typescript. Draft compared: ps4.tex:2652-2724.
% AI-transcribed by Claude Opus 5.5 (claude-opus-5-5, adaptive thinking, effort high), Message Batches API, batch ps2 (msgbatch_01NVwr1Yuo8qJf5rykeGBXPK), 2026-11-12; not yet checked by eye.
\scanpage{134-5/53}
\mspage{240}

The three categories in (12) are endowed respectively with sets of arrows
$W$ , $\underline{W}_A$ , $\underline{W}$ , and the two functors satisfy the conditions
\[ (13) \qquad W = (i^{*})^{-1}(\underline{W}_A) \quad \text{and} \quad \underline{W}_A = (i_A)^{-1}(\underline{W}) \quad , \quad \text{hence} \quad W = (i_M)^{-1}(\underline{W}) \quad , \]
where
\[ i_{M,i}\colon \ M \longrightarrow (\mathrm{Cat}) \]
is the composition $i_A i^{*}$ . Therefore we get a functor
\[ (14) \qquad \overline{i_{M,i}} \ : \ \mathrm{Hot}_M \longrightarrow (\mathrm{Hot}) \quad , \]
which a priori depends upon the choice of $(A,i)$. If we admit strong sa-%
turation of $\underline{W}$, it follows that this functor is ``conservative'', namely
an arrow in $\mathrm{Hot}_M$ is an isomorphism, provided it's image in (Hot) is.

To define (11) in terms of (14), we have to describe merely a tran-%
sitive system of isomorphisms between the functors (14), for varying
pair $(A,i)$. Therefore, consider two such $(A,i)$ and $(A',i')$, choose
a small \inserted{full} subcategory \inserted{B} of $M_{as}$ containing both $i(A)$ and $i'(A')$ (therefore
$B$ is a generating subcategory for the asphericity structure $M_{as}$ ), and
consider the \unsure{inclusion} functor $j\colon B \hookrightarrow M$. From prop.1 b) it follows
that the functors
\[ f\colon A \longrightarrow B \quad , \quad f' : A' \longrightarrow B \]
induced by $i,i'$ are aspheric. Using this, and the criterion (i) of
asphericity (prop. 1 on p. 214) we immediately get two isomorphisms
\diagram{\begin{tikzcd}[column sep=small]
 & \overline{i_{M,\unsure{\over{i}{j}}}} \ar[dl, dash, "\hand{\simeq}"'] \ar[dr, dash, "\hand{\simeq}"] & \\
\overline{i_{M,\unsure{i}}} & & \overline{i_{M,i'}}
\end{tikzcd}}
,

and hence an isomorphism
\[ \overline{i_{M,i}} \xrightarrow{\hand{\sim}} \overline{i_{M,i'}} \quad , \]
depending only on the choice of $B$. As a matter of fact, it does'nt
depend on this choice\over{,}{.} \hand{\unsure{$\leftarrow$}} \over{t}{T}o see this, we are reduced to comparing
the two \unsure{isomorphisms} arising from a $B$ and a $B'$ such that $B \subset B'$, in \xout{whi}
which case the inclusion functor $g\colon B \to B'$ is aspheric\inserted{(prop 1)}. We leave the
detail\hand{s} of the verification\unsure{\deleted{,}} \hand{(}consisting mainly of some diagram chasing
and some compatibilities between the $\hand{\lambda}_i(F)$ -isomorphisms of prop.1 on
page 214\hand{)} to the sceptical reader. As for transitivity for a triple
$(A,i),(A',i'),(A'',i'')$, it \inserted{now} follows at once, by using a $B$ suitable \deleted{\unsure{sim}} \inserted{simul-}

\margin{53}
